\documentclass[10pt,a4paper]{amsart}
\usepackage[margin=19mm]{geometry}
\usepackage{amsmath,amssymb,mathtools}
\usepackage[scr=rsfs,cal=euler]{mathalfa}
\usepackage{xfrac}
\usepackage[unicode,hidelinks,bookmarks=false]{hyperref}
\newcommand{\ie}{i.e.\ }
\newcommand{\cf}{cf.\ }
\newcommand{\resp}{respectively\ }
\newcommand{\Subsection}{\subsection}
\providecommand{\nobreakdash}{\nobreak\hbox{-}}
\setlength{\emergencystretch}{2em}
\multlinegap0pt
\usepackage{amssymb}
\usepackage{xcolor}
\usepackage{mathtools}
\usepackage[shortlabels]{enumitem}
\usepackage{tikz}
\usepackage{float}
\newtheorem{theorem}{Theorem}
\title{The P-Vertex Problem for Graphs with Perfect Matchings}
\author{G. Arunkumar, U. S. Jerisha$^{\ast}$}
\date{Source: arXiv:2608.02107v1 (CC BY 4.0)}
\begin{document}
\maketitle

\noindent\emph{Source and licence.} The P-Vertex Problem for Graphs with Perfect Matchings. Version arXiv:2608.02107v1 (CC BY 4.0), distributed by arXiv under the Creative Commons Attribution 4.0 International licence, http://creativecommons.org/licenses/by/4.0/, as recorded in the arXiv licence metadata for this exact version and shown on the abstract page https://arxiv.org/abs/2608.02107v1. Full licence notice: \texttt{licenses/arxiv-2608.02107-CC-BY-4.0.txt}.\par\smallskip

\noindent\textit{Editorial scope.} Counted unit: the article's theorem thm:main (source0.tex lines 520--525). The article's complete original proof of it is delivered with it (source0.tex lines 527--1125, one \texttt{proof} environment of 599 lines). No mathematics is written, repaired or re-derived: the statement and the proof are the article's own lines, in the article's own notation, and no retained line is substituted or re-flowed.  No further source-local context is needed: every cross-reference of every retained line resolves inside the two delivered spans.  Nothing was omitted from any retained span, so the package carries no omission record.  No retained line contains a citation, so the wrapper carries no bibliography and no bibliography entry.  No retained line contains a figure environment, an image-inclusion command or drawing code, so no image is delivered and the package declares no asset dependency.  Editorial formatting only, in addition: an AMS \texttt{amsart} preamble that restates the article's own declarations and macros used by the retained lines, namely the environments \texttt{theorem} on separate counters, together with the standard packages those lines need.  The retained environments print in this document as Theorem 1 in order of appearance, whereas the article prints the same items with the numbering of its own full document.  The complete native source of the article as distributed in the arXiv e-print for this exact version is bundled unchanged as \texttt{sources/206670/source0.tex}.
\begin{theorem}\label{thm:main}
Let \(G\) be a graph with no isolated vertices. Then
\(
p(G)\leq 2.
\)
\end{theorem}

\begin{proof}
Let \(M\) be a maximal matching of \(G\). Let \(W\) be the set of
vertices saturated by \(M\) and let
\(
U=V(G)\setminus W.
\)
Since \(M\) is maximal, \(U\) is an independent set. Indeed, if two
vertices of \(U\) were adjacent, then their joining edge could be added
to \(M\), contradicting the maximality of \(M\).

We construct a non-singular matrix in \(S(G)\) for which every vertex
of \(W\) is a \(P\)-vertex and when \(U\neq\varnothing\), a second
non-singular matrix in \(S(G)\) for which every vertex of \(U\) is a
\(P\)-vertex.

Throughout the proof, for distinct vertices \(i\) and \(j\), we write
\(A(i,j)\) for the principal submatrix obtained by deleting the rows
and columns indexed by \(i\) and \(j\).

\medskip

\noindent
\textbf{Step 1: A matrix for the vertices in \(W\).}

Suppose
\(
M=\bigl\{\{x_1,y_1\},\ldots,\{x_t,y_t\}\bigr\}.
\)
Thus,
\(
W=\{x_1,y_1,\ldots,x_t,y_t\}.
\)
Order the vertices so that the vertices in \(W\) occur first, with the
endpoints of each matching edge consecutive, followed by the vertices
in \(U\).

Let
\(
B=
\begin{bmatrix}
0&1\\
1&0
\end{bmatrix}
\)
and define the block diagonal matrix
\(
A_W^0=
\operatorname{diag}
\bigl(
\underbrace{B,\ldots,B}_{t\text{ copies}},
I_{|U|}
\bigr).
\)
Since \(\det B=-1\), the matrix \(A_W^0\) is non-singular.

For each \(v\in W\), let \(d_v\) be a variable corresponding to the
\(v^\text{th}\) diagonal entry, and write
\[
\mathbf d=(d_v)_{v\in W}\in\mathbb R^{|W|}.
\]

Let \(\epsilon\in\mathbb R\). Define \(A_W(\mathbf d,\epsilon)\) to be
the symmetric matrix whose entries are given as follows:

\[
\bigl(A_W(\mathbf d,\epsilon)\bigr)_{vv}
=
\begin{cases}
d_v, & v\in W,\\
1,   & v\in U.
\end{cases}
\]
and for distinct vertices \(v,z\),
\[
\bigl(A_W(\mathbf d,\epsilon)\bigr)_{vz}
=
\begin{cases}
1,
& \{v,z\}\in M,\\
\epsilon,
& \{v,z\}\in E(G)\setminus M,\\
0,
& \{v,z\}\notin E(G).
\end{cases}
\]

Let \(\mathbf d^0=\mathbf 0\). Then
\(
A_W(\mathbf d^0,0)=A_W^0.
\)
Define
\(
F_W:\mathbb R^{|W|}\times\mathbb R
\longrightarrow
\mathbb R^{|W|}
\)
by
\(
F_W(\mathbf d,\epsilon)
=
\bigl(
\det A_W(\mathbf d,\epsilon)(v)
\bigr)_{v\in W}.
\)
Since the entries of \(A_W(\mathbf d,\epsilon)\) depend polynomially
on \((\mathbf d,\epsilon)\), the map \(F_W\) is continuously
differentiable.

For every \(v\in W\), deleting the row and column indexed by \(v\)
from \(A_W^0\) leaves a \(1\times1\) zero block arising from the
matching edge containing \(v\). Therefore,
\(
F_W(\mathbf d^0,0)=\mathbf 0.
\)

We now compute the Jacobian of \(F_W\) with respect to \(\mathbf d\) at the point \((\mathbf{d}^0,0)\).

For \(v,z\in W\),
\[
\left.
\frac{\partial}{\partial d_z}
\det A_W(\mathbf d,\epsilon)(v)
\right|_{(\mathbf d^0,0)}
=
\begin{cases}
0, & v=z,\\[1mm]
\det A_W^0(v,z), & v\neq z.
\end{cases}
\]
The first equality holds because the entry \(d_v\) does not occur in
\(A_W(\mathbf d,\epsilon)(v)\). 
%The second follows from differentiation with respect to a diagonal entry \textcolor{purple}{(This is not the explanation for this)}.
The second equality follows from the calculation that if \(v\neq z\), then \(d_z\) appears only as the \((z,z)\)-entry of
\(A_W(\mathbf{d},\epsilon)(v)\). In the cofactor expansion of
\(\det\bigl(A_W(\mathbf{d},\epsilon)(v)\bigr)\) along the row corresponding to \(z\), the term
corresponding to the \((z,z)\)-entry contributes
\(
d_z\cdot\det\bigl(A_W(\mathbf{d},\epsilon)(v,z)\bigr),
\)
whereas all the other terms involve entries of that row outside the
\(z^{\text{th}}\) column and hence do not contain \(d_z\).

If \(v\) and \(z\) are the endpoints of the same edge of \(M\), then
deleting \(v\) and \(z\) removes one entire copy of \(B\). Hence,
\[
\det A_W^0(v,z)=(-1)^{t-1}\neq0.
\]

If \(v\) and \(z\) belong to two different matching edges, then
\(A_W^0(v,z)\) contains two \(1\times1\) zero diagonal blocks arising from two different matching edges containing $v$ and $z$  and
therefore
\(
\det A_W^0(v,z)=0.
\)

Consequently, with respect to the ordering
\(
x_1,y_1,x_2,y_2,\ldots,x_t,y_t,
\)
the Jacobian is a block diagonal matrix given by,
\[
D_{\mathbf d}F_W(\mathbf d^0,0)
=
\operatorname{diag}
\left(
\underbrace{
\begin{bmatrix}
0&c_W\\
c_W&0
\end{bmatrix},
\ldots,
\begin{bmatrix}
0&c_W\\
c_W&0
\end{bmatrix}
}_{t\text{ copies}}
\right),
\]
where
\(
c_W=(-1)^{t-1}.
\)
Every block is non-singular and hence
\(
D_{\mathbf d}F_W(\mathbf d^0,0)
\)
is non-singular.

By the Implicit Function Theorem, there exist \(\delta_W>0\) and a unique
continuously differentiable function
\(
\varphi_W:(-\delta_W,\delta_W)\longrightarrow\mathbb R^{|W|}
\)
such that
\(
\varphi_W(0)=\mathbf d^0
\)
and
\(
F_W(\varphi_W(\epsilon),\epsilon)=\mathbf 0
\)
for every \(\epsilon\in(-\delta_W,\delta_W)\).
Since
\(
\det A_W(\mathbf d^0,0)=\det A_W^0\neq0,
\)
continuity of the determinant allows us, after decreasing \(\delta_W\)
if necessary, to assume that
\(
\det A_W(\varphi_W(\epsilon),\epsilon)\neq0
\)
for every \(\epsilon\in(-\delta_W,\delta_W)\).

Choose
\(
\epsilon_W\in(-\delta_W,\delta_W)\setminus\{0\}
\)
and set
\(
\mathcal A_W
=
A_W(\varphi_W(\epsilon_W),\epsilon_W).
\)
Every edge of \(M\) has matrix entry \(1\), every edge of
\(E(G)\setminus M\) has matrix entry \(\epsilon_W\neq0\), and every
non-edge has matrix entry zero. Hence,
\(
\mathcal A_W\in S(G).
\)
Moreover,
\(
\det\mathcal A_W\neq0
\)
and, for every \(v\in W\),
\(
\det\mathcal A_W(v)=0.
\)

Thus,
\(
m_{\mathcal A_W}(0)=0
\)
and
\(
m_{\mathcal A_W(v)}(0)\geq1.
\)
By the Cauchy interlacing theorem,
\[
m_{\mathcal A_W(v)}(0)-m_{\mathcal A_W}(0)\leq1.
\]
It follows that
\(
m_{\mathcal A_W(v)}(0)=1
\)
and therefore every \(v\in W\) is a \(P\)-vertex of
\(\mathcal A_W\).

If \(U=\varnothing\), then \(W=V(G)\), so \(\mathcal A_W\) already
shows that \(p(G)=1\). We may therefore assume that
\(U\neq\varnothing\).

\medskip

\noindent
\textbf{Step 2: A matrix for the vertices in \(U\).}

Since \(G\) has no isolated vertices and \(U\) is independent, every
vertex of \(U\) has a neighbour in \(W\). For each \(u\in U\), choose
one such neighbour and denote it by \(f(u)\).

For \(w\in W\), define
\(
U_w=\{u\in U:f(u)=w\}.
\)
Let
\(
W'=\{w\in W:U_w\neq\varnothing\}
=
\{w_1,\ldots,w_q\}.
\)
For each \(1\leq i\leq q\), write
\(
U_i=U_{w_i}
=
\{u_{i,1},\ldots,u_{i,n_i}\},
\
n_i=|U_i|\geq1.
\)
The sets \(U_1,\ldots,U_q\) form a partition of \(U\).

For each \(i\), define the \((n_i+1)\times(n_i+1)\) matrix
\[
S_i=
\begin{bmatrix}
n_i-1&\mathbf 1_{n_i}^{T}\\
\mathbf 1_{n_i}&I_{n_i}
\end{bmatrix},
\]
where the first row and column correspond to \(w_i\), and the
remaining rows and columns correspond to the vertices in \(U_i\).

By the Schur complement formula,
\[
\det S_i
=
\det(I_{n_i})
\left(
n_i-1-\mathbf 1_{n_i}^{T}\mathbf 1_{n_i}
\right)
=(n_i-1)-n_i
=-1.
\]
Thus every \(S_i\) is non-singular.

If \(u\in U_i\), then deleting the row and column corresponding to 
\(u\) from the matrix \(S_i\) gives,
\[
\det S_i(u)
=
(n_i-1)-(n_i-1)
=0.
\]

If \(n_i\geq2\) and \(u,v\in U_i\) are distinct, then
\[
\det S_i(u,v)
=
(n_i-1)-(n_i-2)
=1.
\]

With respect to the ordering
\(
w_1,U_1,w_2,U_2,\ldots,w_q,U_q,W\setminus W',
\)
define
\[
A_U^0
=
\operatorname{diag}
\left(
S_1,\ldots,S_q,I_{|W\setminus W'|}
\right).
\]

Since every diagonal block is non-singular,
\(
\det A_U^0\neq0.
\)
Moreover, for every \(u\in U\),
\(
\det A_U^0(u)=0.
\)

We now choose the diagonal variables used in the Implicit Function
Theorem. Define a map
\(
\psi:U\longrightarrow V(G)
\)
as follows. If \(u\in U_i\), set
\(
\psi(u)=
\begin{cases}
w_i, & n_i=1,\\
u,   & n_i\geq2
\end{cases}
\).
The map \(\psi\) is injective. 

Thus, we may assign one independent
diagonal variable \(h_u\) to the diagonal entry indexed by
\(\psi(u)\), for each \(u\in U\). Write
\(
\mathbf h=(h_u)_{u\in U}\in\mathbb R^{|U|}.
\)
Let \(\mathbf h^0\) denote the values of these diagonal entries in
\(A_U^0\). Explicitly,
\[
h_u^0=
\begin{cases}
0, & u\in U_i\text{ and }n_i=1,\\
1, & u\in U_i\text{ and }n_i\geq2
\end{cases}.
\]

Let
\(
E_0=\bigl\{\{w_i,u\}:1\leq i\leq q,\ u\in U_i\bigr\}.
\)
Define \(A_U(\mathbf h,\epsilon)\) as follows. Its diagonal entries
indexed by \(\psi(u)\), \(u\in U\), are replaced by the corresponding
variables \(h_u\), while all other diagonal entries remain equal to
their values in \(A_U^0\). For distinct vertices \(v,z\), define
\[
\bigl(A_U(\mathbf h,\epsilon)\bigr)_{vz}
=
\begin{cases}
1,
& \{v,z\}\in E_0,\\
\epsilon,
& \{v,z\}\in E(G)\setminus E_0,\\
0,
& \{v,z\}\notin E(G).
\end{cases}
\]

Then
\(
A_U(\mathbf h^0,0)=A_U^0.
\)
Define
\(
F_U:\mathbb R^{|U|}\times\mathbb R
\longrightarrow
\mathbb R^{|U|}
\)
by
\[
F_U(\mathbf h,\epsilon)
=
\bigl(
\det A_U(\mathbf h,\epsilon)(u)
\bigr)_{u\in U}
\]
This is a continuously differentiable map and
\(
F_U(\mathbf h^0,0)=\mathbf 0.
\)

We compute the Jacobian of \(F_U\) with respect to \(\mathbf h\).
For \(u,v\in U\),
\[
\left.
\frac{\partial}{\partial h_v}
\det A_U(\mathbf h,\epsilon)(u)
\right|_{(\mathbf h^0,0)}
=
\begin{cases}
0,
& \psi(v)=u,\\[1mm]
\det A_U^0(u,\psi(v)),
& \psi(v)\neq u.
\end{cases}
\]

Let
\(
c_U=(-1)^{q-1}.
\)
We examine the Jacobian block corresponding to each set \(U_i\).

Suppose first that \(n_i=1\), and write \(U_i=\{u\}\). Then
\(
\psi(u)=w_i,
\)
and therefore
\[
\left.
\frac{\partial}{\partial h_u}
\det A_U(\mathbf h,\epsilon)(u)
\right|_{(\mathbf h^0,0)}
=
\det A_U^0(u,w_i).
\]

Deleting both \(u\) and \(w_i\) removes the entire block \(S_i\).
All the remaining blocks are non-singular and hence
\[
\det A_U^0(u,w_i)=c_U\neq0.
\]
Thus, the Jacobian block corresponding to \(U_i\) is the
\(1\times1\) matrix
\(
[c_U].
\)

Now suppose that \(n_i\geq2\). For \(u,v\in U_i\), we have
\(\psi(v)=v\). Therefore,
\[
\left.
\frac{\partial}{\partial h_v}
\det A_U(\mathbf h,\epsilon)(u)
\right|_{(\mathbf h^0,0)}
=
\begin{cases}
0, & u=v,\\
c_U, & u\neq v.
\end{cases}
\]
Indeed, when \(u\neq v\), deleting the two leaves \(u\) and \(v\)
from \(S_i\) gives a block of determinant \(1\), while every other
block contributes its determinant.

If \(u\in U_i\) and \(v\in U_j\) with \(i\neq j\), then
\(
\det A_U^0(u,\psi(v))=0.
\)
To see this, observe that deleting \(u\) makes the block \(S_i\)
singular, and deleting \(\psi(v)\), which belongs to a different
block, does not remove this singularity.

Consequently, after ordering the vertices of \(U\) according to the
partition
\(
U=U_1\mathbin{\dot\cup}\cdots\mathbin{\dot\cup}U_q,
\)
the Jacobian is a block diagonal matrix given by,
\[
D_{\mathbf h}F_U(\mathbf h^0,0)
=
\operatorname{diag}(J_1,\ldots,J_q),
\]
where
\[
J_i=
\begin{cases}
[c_U], & n_i=1,\\[2mm]
c_U\bigl(\mathbf J_{n_i}-I_{n_i}\bigr), & n_i\geq2,
\end{cases}
\]
and \(\mathbf J_{n_i}\) denotes the \(n_i\times n_i\) all-one
matrix.

For \(n_i\geq2\), the eigenvalues of
\(\mathbf J_{n_i}-I_{n_i}\) are
\(
n_i-1
\text{ and } 
-1
\)
with multiplicities \(1\) and \(n_i-1\), respectively. Hence every
\(J_i\) is non-singular. Therefore,
\(
D_{\mathbf h}F_U(\mathbf h^0,0)
\)
is non-singular.

By the Implicit Function Theorem, there exist \(\delta_U>0\) and a unique
continuously differentiable function
\(
\varphi_U:(-\delta_U,\delta_U)\longrightarrow\mathbb R^{|U|}
\)
such that
\(
\varphi_U(0)=\mathbf h^0
\)
and
\(
F_U(\varphi_U(\epsilon),\epsilon)=\mathbf 0
\)
for every \(\epsilon\in(-\delta_U,\delta_U)\).
Since
\(\det A_U(\mathbf h^0,0)=\det A_U^0\neq0,
\)
we may decrease \(\delta_U\), if necessary, so that
\(
\det A_U(\varphi_U(\epsilon),\epsilon)\neq0
\)
for every \(\epsilon\in(-\delta_U,\delta_U)\).
Choose
\(
\epsilon_U\in(-\delta_U,\delta_U)\setminus\{0\}
\)
and set
\(
\widetilde{\mathcal{A}}_U
=
A_U(\varphi_U(\epsilon_U),\epsilon_U).
\)
By construction,
\(
\widetilde{\mathcal{A}}_U\in S(G)
\text{ and }
\det\widetilde{\mathcal{A}}_U\neq0.
\)
Furthermore, for every \(u\in U\),
\(
\det\widetilde{\mathcal{A}}_U(u)=0.
\)

Since the above construction is carried out with respect to a temporary ordering of the vertices, the matrix obtained need not correspond to the original labeling of \(G\) used in Step 1. Let \(P\) be the permutation matrix that restores the original ordering in Step 1 and define
\(
\mathcal{A}_U=P^{T}\widetilde{\mathcal{A}}_U P.
\)
Since simultaneous permutation of rows and columns preserves symmetry, non-singularity, principal minors, and the graph of a matrix, it follows that \(\mathcal{A}_U\in S(G)\), \(\det(\mathcal{A}_U)\neq0\), and \(\det(\mathcal{A}_U(u))=0\) for every \(u\in U\).
As in Step 1, Cauchy interlacing theorem implies that
\(
m_{\widetilde{\mathcal{A}}_U(u)}(0)=1
\)
for every \(u\in U\). Thus, every vertex of \(U\) is a \(P\)-vertex
of \(\mathcal{A}_U\).

The matrix \(\mathcal A_W\) covers all vertices of \(W\) makes all the vertices of $W$ as $P$-vertices and the
matrix \(\mathcal A_U\) covers all vertices of \(U\) makes all the vertices of $U$ as $P$-vertices. Therefore,
\(
p(G)\leq2.
\) 

\end{proof}

\end{document}
